%! A two-column LaTeX fixture for scroll correspondence.
%!
%! The case that matters: one source line can own regions in *both* columns, so a
%! correspondence rule that picks the first region snaps to the left column while
%! the reader is scrolling down the right one. `multicols` is used rather than the
%! `twocolumn` class option because it keeps the source linear — the reader's
%! document is one column of text that happens to be typeset in two.
%!
%! **No `lipsum`.** It needs a package fetch, and the test profile's cache does not
%! have it, so the fixture failed to compile with `File 'lipsum.ltd' not found`.
%! Plain repeated prose fills two columns just as well and needs nothing beyond
%! `multicol`, which every TeX Live scheme carries.
\documentclass[11pt]{article}
\usepackage{multicol}
\title{Two Columns}
\author{MDraft}
\begin{document}
\maketitle

\section{Left and right}\label{sec:columns}

\begin{multicols}{2}
The first paragraph begins in the left column and, if it is long enough, will
continue into the right one. That is exactly the case the correspondence has to
get right: the same source line owns regions on both sides of the page, and a
rule that always takes the first region would jump back to the left column on
every scroll event.

A second paragraph continues the argument. The reader scrolling through this
document expects the preview to move downward continuously, not to oscillate
between the columns. Each of these paragraphs is long enough to wrap several
times, so the column break falls somewhere inside one of them rather than neatly
between two.

A third paragraph, so there is enough text to fill both columns and to scroll
past the first page. The correspondence is only interesting when there is more
content than fits on one screen, and a two-column layout reaches that point
sooner than a single-column one.

A fourth paragraph adds more lines. Without it the environment would end before
the right column filled, and the case this fixture exists for would not arise.

A fifth paragraph, to be certain. The preview should keep moving down as the
editor scrolls, whichever column the visible line happens to be typeset in.

A sixth paragraph closes the environment. Text after it maps to a single-column
region, which is the contrast the rule has to handle.
\end{multicols}

\section{After the columns}

Text after the environment, which must map back to a single-column region. It
exists so the fixture covers both layouts in one document.

\end{document}
